\section{An approximation estimate at every smaller scale}
\label{sec:compactness}

The electrical estimates of Theorem~\ref{thm:network-bounds} control
each fixed annulus after a threshold which may depend on that annulus
and on the sample.  Spectral compactness requires one estimate valid
at all remaining discrete scales.  We obtain it from the address
estimates and finite contour views of~\cite{P}.  The new ingredients
are a rank estimate with one marked time, a bound on the number of
widely separated times reached by short contour explorations, and a
cover by circuits which avoids the cost of reaching individual
vertices.

For a set \(W\) of original edges, counted with their map
multiplicities, write
\[
 \mathcal E_m(f;W)=\sum_{e\in W}\abs{f(e^+)-f(e^-)}^2.
\]
The choice of orientation \(e^-,e^+\) is immaterial, and a loop
contributes zero.  On a duration-one sphere \(m=2n\), this is the
restriction of \(\mathcal E_n^0\) to \(W\).

An interval used below is a set \(I_m\) of consecutive letter cuts
inside one used block of a protected chart.  Its endpoints are
specified relative to that block, its length divided by \(m\) has a
positive finite limit, and its limiting tour image has compact
clearance inside a larger protected neighborhood.  The chart records
a graph \(G_m^+\) which contains the interval's vertices and all
edges and complete stars in an intermediate neighborhood, and is
itself reconstructed with room inside the larger neighborhood.
These are entries in the countable atlas of
Proposition~\ref{prop:readout-extension}; activation refers to its
strict continuum conditions.

Split \(I_m\) into nested dyadic halves, rounding integer sizes at
each split.  Retain singleton and empty blocks at later levels.  For
a nonempty level-\(j\) block \(I\), define
\[
 (P_j^{I_m}f)(i)=\frac1{\abs I}\sum_{k\in I}f(v(k)),
 \qquad i\in I.
\]
Thus \(P_j^{I_m}\) is the orthogonal projection of the cut-value
function onto level-\(j\) block constants for counting measure; it
is the identity after every block is a singleton.

\begin{theorem}[All-level block approximation]
\label{thm:block-approximation}
From an arbitrary sequence of finite spheres, take a further joint
representation extraction as in Proposition~\ref{prop:geometry}
and Proposition~\ref{prop:readout-extension}, retaining the interval
readouts defined below.  Put \(m=2n\) and
\(p=7/10\).  Simultaneously for the countable protected interval
atlas, the following holds almost surely on activation of an entry
\((I_m,G_m^+)\).  There are finite sample-dependent constants
\(j_*\), \(C\), and \(n_*\) such that, at every index \(n\ge n_*\)
in the extraction, for every function on \(V(G_m^+)\) and every
integer \(j\ge j_*\),
\begin{equation}
\label{eq:block-approximation}
 \frac1m\sum_{i\in I_m}
   \abs{f(v(i))-(P_j^{I_m}f)(i)}^2
 \le C\,2^{-pj} \mathcal E_n^0(f;G_m^+).
\end{equation}
This includes all levels down to and beyond singleton cuts and holds
for all functions on the same event.

In the two-rooting extraction of Proposition~\ref{prop:geometry},
one can choose finitely many activated entries from the two rooted
atlases such that every corner occurrence belongs to at least one
of their intervals, in its corresponding rooted order, for all
sufficiently large \(n\).
\end{theorem}

We first work with the bilateral word and prove a tight bound on the
smallest sufficient \(j_*\).  Only after that probability estimate
will we regard the threshold as a local readout and transfer
almost-sure finiteness of its extracted limit on activation to sphere charts.

\subsection{Selected words and the address estimates}

Let \(\PP_m^W\) be the product law on bilateral inventory words.
The two burger types, their two rigid orders, and the flexible
order have probabilities
\[
 \frac14,\quad\frac14,\quad
 \frac{2-\sqrt2}{4},\quad\frac{2-\sqrt2}{4},\quad
 \frac{\sqrt2-1}{2},
\]
respectively.  An order consumes the freshest allowed burger.
Resolve flexible orders using the bilateral matching.  The subscript
records the time unit \(m\), not a change of the iid letter law.
The two resolved scalar contours are \(Z^1,Z^2\); their increments
are not asserted independent.  These are the FK-Ising probabilities
of \cite[Proposition~3.1]{C}.  Work in an ordinary bilateral contour
coupling with its reference tour projection; extra field coordinates
use the attachment convention of Proposition~\ref{prop:readout-extension}.
We use a
\emph{scalar link} only for a pair of integer cuts \(s,t\) satisfying
\[
 Z^c(s)=Z^c(t)=\min_{[s\wedge t,s\vee t]}Z^c
 \quad\text{for one }c\in\{1,2\}.
\]
Links can have arbitrary chronological duration.  A time jump is
an arbitrary pair of cuts whose distance is bounded as specified
below.  Put
\[
 r_m(x)=m e^{-x},\qquad
 \vartheta=\frac12+\frac2{\gamma^2}=\frac76.
\]
Every sequential statement in this subsection has
\begin{equation}
\label{eq:selected-units}
 N=N_m\longrightarrow\infty,\qquad
 r_m(bN)\longrightarrow\infty
\end{equation}
for the fixed upper log-range constant \(b>1\).

We record the selection protocol because it will also allow a union
over many levels.  Fix a compact starting window in time units, a
spatial compactness and positive-clearance cutoff, and
\(\sigma_m\downarrow0\) bounding every used seed unit divided by \(m\).
For any fixed
loss \(\epsilon>0\), the preselection construction of
\cite[Section~11]{P} supplies a master event \(\mathcal T_m\),
deterministic enclosures \([-C_d m,C_d m]\), and a deterministic
\(D_m^{\max}\to\infty\).  Its failure probability on the fixed
cutoff is at most \(2\epsilon+o(1)\).  On \(\mathcal T_m\),
simultaneously for \(d\le D_m^{\max}\), every exploration of at most
\(d\) links or jumps of size at most \(dm\) from the starting
window remains in the \(d\)th enclosure.  If its jumps have size
at most \(d\sigma_m m\), all its projected points are within a
deterministic \(\chi_m\downarrow0\) of its base.  This includes
first- and last-return witnesses.  For each fixed \(d\), the
bilateral protected-window result \cite[Proposition~4.9]{C},
after a deterministic compact cutoff, gives the vanishing flag
mesh and exact complete stars needed for the fixed intermediate
neighborhood on that enclosure.  Include its eventual conclusion
with a further prescribed summable loss before selection.
Slowing \(D_m^{\max}\) makes these conclusions simultaneous as well.

A selected law \(\mathbb Q_m\) in this section is supported on
\(\mathcal T_m\) and the stated clearance cutoff.  Its raw-word
marginal satisfies
\begin{equation}
\label{eq:selected-density}
 \frac{d\mathbb Q_m^W}{d\PP_m^W}\le e^{\kappa_m},
 \qquad \kappa_m=o(N).
\end{equation}
Finite lists may be randomized functions of the word and of the
coupled environment.  To apply the companion on its augmented
probability space, give the reference word the selected law's
conditional selector kernel given that word.  This kernel can be
realized by an independent uniform variable; the augmented reference
still has word marginal \(\PP_m^W\), and the full augmented density
is the density in~\eqref{eq:selected-density}.  The selected environment
and bases precede the independent viewing level or seed mesh.
Seed-specific lists may then use their seed levels; their complete
union of positions and access witnesses precedes the independent
chronological address grid.  Subsequent
finite-list cutoffs may grow only along a slower diagonal below
\(D_m^{\max}\).  These are the hypotheses on every selected law
below; we will not transfer a fresh high-probability master event
through~\eqref{eq:selected-density}.

Here is the quantitative address input.  In an enclosure of length
\(h_m^{\mathrm{enc}}=c_{\mathrm{enc}}m\), fix an integer \(R\ge8\), put
\(\ell_R=\log R\), and use generations
\(0\le k\le J_R=\lfloor bN/\ell_R\rfloor\) of cells of width
\(\delta_k=h_m^{\mathrm{enc}}R^{-k}\).  There are \(h_{\mathrm b}\)
distinguished bases marked through \(J_R\).  Use distinct probe
generations.  At each probe generation,
every endpoint and intermediate point in each access chain from
a distinguished base is
marked through that generation; its links are scalar links or
jumps of size at most \(K\delta_k\).  The fixed \(K\) also bounds
the marker and edge entries in every probe list.  Include all distinguished
base entries in the list at terminal generation \(J_R\), merging
them into an existing list there or adding that generation as one
probe.  Enlarge \(K\) to bound the marker and edge entries in this
augmented list as well.  Let \(P_N\) count the resulting distinct
probes and \(M_{\rm term}\) count all marker entries, so
\(M_{\rm term}\le KP_N\).  The count below fixes these augmented
probe depths; in our applications they are determined by an
independent mesh phase held as side information.  The entire
augmented marker system is chosen before the independent
\(\Xi\sim\operatorname{Unif}[0,h_m^{\mathrm{enc}})\) translates the nested grids.

At generation \(k\), let \(m_k\) be the number of occupied cells,
\(b_k\) the number containing a base, and \(D_k\) the number of jump
entries split between distinct cells.  For each color, project every
active recorded chord of that color to the cells containing its endpoints and take
all occupied cells, including isolated ones, as vertices; these graphs
exclude jumps, as in \cite[Definition~10.6 and Lemma~10.7]{P}.
Define graph rank as the number of vertices minus the number of
connected components, so projected loops and repeated edges add no
rank, and let \(e_k\) be the sum of the two color ranks.  Write
\[
 T_{\rm occ}=\sum_k m_k,\qquad R_{\rm rk}=\sum_k e_k.
\]

\begin{proposition}[Address input from the companion]
\label{prop:address-input}
For the marker systems just described, with fixed \(b,R,K\), the
following statements hold.
\begin{enumerate}[label=(\roman*)]
\item If \(c_k\) is the number of components of the union of the
two colored graphs, then
\[
 c_k\le b_k+D_k,\qquad e_k\ge m_k-c_k.
\]
Moreover
\[
 \EE_\Xi\sum_kD_k\le C_{K,R}M_{\rm jump}.
\]
If the bases lie in an interval of length \(C r_m(N)\), then
\(\EE_\Xi\sum_{k\le N/\ell_R}(b_k-1)\le C_{C,R}\).
\item For each \(\varepsilon>0\), a fixed sufficiently large \(R\)
gives, with conditional grid probability at least \(7/8\), disjoint
regular sampling slots after deletion of at most
\(C_RM_{\rm term}+\varepsilon T_{\rm occ}\) node occurrences.
For each sufficiently small fixed \(\varepsilon>0\), \(R\) and
the deletion tolerances may be fixed so that every such code
\(\mathfrak c\) obeys
\begin{equation}
\label{eq:fixed-code-charge}
 \PP_m^W\{\mathfrak c\text{ is realized}\}
 \le
 \exp\{-\vartheta\ell_R R_{\rm rk}
       +\varepsilon\ell_R T_{\rm occ}
       +C_{R,\varepsilon}M_{\rm term}\}.
\end{equation}
Realization includes all its actual chord constraints.  The bound
is uniform in the fixed grid displacement and in the exact
chord endpoints inside the recorded cells, once the minimum
sampling radius diverges.
\item For each fixed displacement and fixed augmented distinct
probe depths, and every \(0<\eta\le K\), the number of decorated codes with
occupancy \(T_{\rm occ}=t\) is at most the exponential of
\begin{equation}
\label{eq:address-count}
 (\ell_R+C_0+\eta)t
 +K P_N\log(K/\eta)+C P_NK\log(K+1)+O(1).
\end{equation}
If the depth list itself is encoded, add
\(\log\binom{J_R+1}{P_N}\); our independent mesh phases are
instead held as side information.
\item Let \(\Delta(\Xi)\) count generations at which the partition
of the active markers into shifted cells differs from the
grid-independent partition obtained by joining consecutive active
positions with gap \(<\delta_k\).  Then
\[
 \EE_\Xi[\Delta(\Xi)\mid\text{markers}]
 \le \frac{2R}{R-1}M_{\rm term}.
\]
At a generation with equal partitions the entire colored graph
and its base labels agree up to cell names.
\end{enumerate}
For \(P_N=O(N/L)\), take the error tolerances and \(R\) fixed first,
then \(N\) and the minimum units to infinity with \(L\to\infty\)
and \(L=o(N)\).  The terminal, jump, and structural errors above
are \(o(N)\) for every fixed \(K\).  One may subsequently let
\(K=K_m\to\infty\) sufficiently slowly, with
\[
 \frac{K_m\log(K_mL)}{L}\longrightarrow0,\qquad
 \frac{K_m\log(2\widehat C_m)}{L}\longrightarrow0,
\]
where \(2\widehat C_m m\) is a growing enclosure length, and with
the composed exploration depth below \(D_m^{\max}\).
The second condition pays for its extra address ancestors.  The
integrated jump error remains \(o(N)\), and each previously fixed
minimum-unit threshold is retained.
\end{proposition}

\begin{proof}
The projected rank and jump inequalities are the spanning-forest
and displaced-grid estimates in~\cite[Section~10]{P}.
The fixed-code estimate is \cite[Proposition~10.15]{P}; its proof
charges a left and right fresh slot for each monochrome rank and
gives the coefficient \(\vartheta=1/2+2/\gamma^2\).
The decorated address count is \cite[Lemma~10.13]{P}: distinct
depths give \(\sum_p m_{k_p}\le T_{\rm occ}\), and the terminal
base list supplies every base's coarser address.  The regular-slot
deletion bound and grid-partition estimate are the corresponding
lemmas in the same section.  Its sparse-diagonal remark gives the
list and jump order of limits; the extra-ancestor estimate in the
proof of \cite[Lemma~11.6]{P} gives the growing-enclosure condition.
In these formulas a generation has log width
\(\ell_R\); this accounts for every conversion between sums over
generations and log lengths.
\end{proof}

We need a packing consequence also at the order of limits in which
the probe stride is held fixed while \(m\) grows.  For \(k\ge1\),
let \(\mathcal R_{m,k}(x)\) be all times reached from the selected
bases by at most \(k\) scalar links or jumps, each jump of size at
most \(k r_m(x)\).  The whole access chain is included.  Bilateral
recurrence below each fixed scalar height makes this set finite.
It is increasing in \(k\), and
\(\mathcal R_{m,k}(x)\subset\mathcal R_{m,k}(y)\) for \(y\le x\).
Define
\[
 \Pi_{m,k,H}(x)=
 \max\{\abs A:A\subset\mathcal R_{m,k}(x),\
       \abs{s-t}>H r_m(x)\ \text{for }s\ne t\text{ in }A\}.
\]

\begin{lemma}[Truncated packing at a typical level]
\label{lem:truncated-packing}
Use the selected protocol above with a fixed number of bases.
Let \(\mathcal I_N\subset[a_-N,bN]\), with fixed \(a_->1\), be
a deterministic log interval of length at least \(cN>0\), and let \(Y\) be uniform on
\(\mathcal I_N\), independently after the selected environment.
Then there is \(C_{\rm pack}<\infty\), independent of each fixed
\(k,M\), such that
\begin{equation}
\label{eq:truncated-packing}
 \sup_{k,M\ge1}\ 
 \lim_{H\to\infty}\limsup_m
 \EE_{\mathbb Q_m}\min\{\Pi_{m,k,H}(Y),M\}
 \le C_{\rm pack}.
\end{equation}
\end{lemma}

\begin{proof}
Choose fixed address parameters within
Proposition~\ref{prop:address-input} so that
\[
 \delta=\vartheta-1-\varepsilon-\frac{C_0+\eta}{\ell_R}>0.
\]
Fix \(k,M\) and a log stride \(L>2\ell_R\).  Sample an independent
translation of a stride-\(L\) mesh in \(\mathcal I_N\).  At each
seed choose a maximum-cardinality \(H\)-packing truncated at \(M\), with all
access witnesses.
Use a fixed measurable ordering of the finite reachable set, and
keep the independent mesh phase as side information in the code count.
Round a seed \(y\) up to generation \(\lceil y/\ell_R\rceil\),
discarding the final incomplete generation if needed.  The seed
depths are distinct because \(L>2\ell_R\).  The terminal base-list
augmentation adds at most one probe and \(h_{\mathrm b}\) fixed
entries, so \(P_N=O(N/L)\) and \(K=O(kM+h_{\mathrm b})\), with
constants allowed to depend on the fixed enclosing depth.
Only then sample the chronological address-grid displacement.
These lists have exactly the order and the form in
Proposition~\ref{prop:address-input}.

We first bound their occupation for this \emph{fixed} \(L\).
Markov's inequality in part (i) gives a set of grid displacements
of fixed positive conditional probability on which
\(\sum D_k\le C_KN/L+O(1)\).  Intersect it with the regular-slot
event, retaining a fixed positive probability.  Use
\(e_k\ge m_k-h_{\mathrm b}-D_k\) in
\eqref{eq:fixed-code-charge} and sum the codes by
\eqref{eq:address-count}.
All terminal and decoration terms at fixed \(K,L\) are
\(O_K(N/L)+O(1)\).  Let \(\mathsf F_t(\Xi)\) be the existence
of a realized regular code with this jump bound and occupancy \(t\).
Uniformly in the displacement,
\begin{equation}
\label{eq:fixed-stride-occupation}
 \log\PP_m^W(\mathsf F_t(\Xi))
 \le -\delta\ell_Rt+
       \bigl[\vartheta h_{\mathrm b}b+O_K(L^{-1})\bigr]N+O(1).
\end{equation}
The constants in the leading term are independent of \(K\).
This calculation retains the terminal errors explicitly; the
asymptotic occupancy statement with \(L=L_m\to\infty\) alone would
not justify the present order of limits.

For two displacements, each cell of one grid meets at most two
cells of the other, hence \(T_{\rm occ}(\Xi)\le
2T_{\rm occ}(\Xi')\).  If the maximum occupancy over displacements
is large, it is at least half as large on every displacement,
including the fixed positive fraction just used.  Integrate
\eqref{eq:fixed-stride-occupation} over the grid and sum over \(t\).
Then apply~\eqref{eq:selected-density}.  It follows that, for a
constant \(C_1\) independent of \(K\), the total log occupation
is at most
\[
 [C_1+O_K(L^{-1})]N
\]
with \(\mathbb Q_m\)-probability tending to one.  The same bound
applies to the occupation restricted to \(\mathcal I_N\).

Here is the conversion to packing, which is the mechanism of
\cite[Lemma~11.4]{P}.  A list chosen at a seed \(y\) remains
reachable throughout the coarser interval \([y-L,y]\).
Use generation \(\lceil x/\ell_R\rceil\) boxes on each continuous
log cell, trimming incomplete endpoint cells.  They are active
throughout \([y-L,y]\) for a seed rounded as above, and their
integrated occupation is at most \(\ell_R T_{\rm occ}\).
If these boxes at level \(x\) have diameter at most
\(c_{\rm box}r_m(x)\), then its chosen terminals stay in distinct
boxes throughout that interval whenever
\(H e^{-L}>2c_{\rm box}\).  Thus the truncated seed count
contributes at least \(L\) times that count to integrated
occupation, apart from the two end intervals and a fixed-width
endpoint strip.  Averaging the
independent mesh translation turns the seed sum into the uniform
integral on \(\mathcal I_N\).  On the occupation-exception event,
bound the truncated packing by \(M\).  Dividing by
\(\abs{\mathcal I_N}\), the end error is at most
\(O(M(L+\ell_R)/\abs{\mathcal I_N})\).  First let \(m\to\infty\) for fixed
\(H,L,k,M\); then take \(H\to\infty\) and
\(L=L(H)\to\infty\) slowly with \(He^{-L}>2c_{\rm box}\).
The term \(O_K(L^{-1})\) vanishes for each fixed \(k,M\), and
the limiting bound is independent of them.  This proves
\eqref{eq:truncated-packing}.
\end{proof}

\begin{lemma}[Few separated times at bounded reach]
\label{lem:reachable-packing}
Fix \(a_0>1\), \(c_0,C_0,\zeta>0\), and a finite integer \(d\).
Under the selected protocol with \(r_m(a_0N)\to\infty\), there
is no sequence of selected laws supported on the existence of
more than \(e^{\zeta N}\) times in the starting window, pairwise
separated by at least \(c_0r_m(N)\), with one common member from
which every member is reached by at most \(d\) scalar links or jumps of size at most
\(C_0r_m(a_0N)\).
\end{lemma}

\begin{proof}
Choose fixed \(1<\alpha<\beta<a_0\) and apply
Lemma~\ref{lem:truncated-packing} with one selected base and
\(\mathcal I_N=[\alpha N,\beta N]\).  For every \(x\in\mathcal I_N\),
the proposed jumps divided by \(r_m(x)\) are at most
\(C_0e^{-(a_0-\beta)N}\), while the proposed separation divided
by \(r_m(x)\) is at least \(c_0e^{(\alpha-1)N}\).  For a fixed
enlargement of \(d\), the proposed list is therefore contained
in \(\mathcal R_{m,k}(x)\) and is an \(H\)-packing for every
fixed \(H\), eventually uniformly in \(x\).  Select the base
and only \(M\) terminals with their witnesses before the grids.
On the asserted event the truncated packing equals \(M\) for
every fixed \(H,M\) and all large \(m\).  Taking the limits in
\eqref{eq:truncated-packing} and then choosing
\(M>C_{\rm pack}\) is a contradiction.
\end{proof}

The other new use of the address budget replaces many nearby bases
by one base and a wider range of smaller scales.  For a connected
component \(C\) of the union of two colored graphs, the sum of the
ranks of their restrictions to \(C\) is at least \(\abs C-1\).
Call it \emph{acyclic} when equality
holds and say it has a \emph{rank cycle} when the rank is at least
\(\abs C\).

\begin{lemma}[One-base rank credit on a wide range]
\label{lem:one-base-credit}
There is a fixed \(B>2\) with the following property.  Use a
sparse address system from Proposition~\ref{prop:address-input}
with one base, upper range \(BN\), and selected
law~\eqref{eq:selected-density}.  Fix \(1<a\le2\) and any
deterministic \(\eta_m\downarrow0\).  The probability that the
projected component containing the base has a rank cycle at all
but \(\eta_mN\) generations with log level
\(k\ell_R\in[aN,BN]\) tends to zero.
\end{lemma}

\begin{proof}
The base is marked at every generation, so \(b_k=1\) and
\(m_k\ge1\).  At a generation where its component has a rank
cycle, that component contributes one rank beyond a union
spanning tree.  Part (i) of Proposition~\ref{prop:address-input}
therefore gives
\[
 e_k\ge m_k-D_k
 \quad\text{there},\qquad
 e_k\ge m_k-1-D_k
 \quad\text{at every other generation}.
\]
On a regular grid with integrated jump error \(o(N)\), the event
in the statement consequently implies
\begin{equation}
\label{eq:one-base-rank}
 \ell_RR_{\rm rk}\ge\ell_RT_{\rm occ}-aN-o(N),
 \qquad
 \ell_RT_{\rm occ}\ge BN+O(1).
\end{equation}
Here the first loss counts the generations before \(aN\) and the
exceptional high generations.  Extra components separated by
jumps have already been charged to \(\sum D_k\).  A permitted
growing enclosure shifts the log origin by
\(O(\log\widehat C_m)=o(N)\); its extra ancestors are included
in the same error.

Fix \(R\) and the errors so that the \(\delta\) in the proof of
Lemma~\ref{lem:truncated-packing} is at least
\((\vartheta-1)/2\).  Combining
\eqref{eq:fixed-code-charge}, \eqref{eq:address-count}, and
\eqref{eq:one-base-rank}, with the sparse errors \(o(N)\), gives
total log weight at most
\[
 -\delta\ell_RT_{\rm occ}+\vartheta aN+o(N)
 \le
 -\frac{\delta}{2}\ell_RT_{\rm occ}
 +\left(\vartheta a-\frac{\delta B}{2}\right)N+o(N).
\]
Choose \(B\) with \(\delta B/2>2\vartheta\) by a fixed margin.
The first negative term makes the sum over \(T_{\rm occ}\)
geometric, and the remaining exponent is \(-cN+o(N)\).
Multiplication by \(e^{\kappa_m}=e^{o(N)}\) preserves convergence
to zero.

This estimate was for regular grids.  The cycle event itself
depends on the grid, so its intersection with a positive fraction
of grids is not automatic.  Since \(M_{\rm term}=o(N)\), part (iv)
provides \(q_m=o(N)\) and \(\alpha_m\to0\) such that
\(\PP_\Xi(\Delta(\Xi)>q_m\mid\text{markers})\le\alpha_m\).
Let \(E^*\) be the corresponding cycle event for the
grid-independent partition, allowing \(q_m\) additional
exceptions.  If \(E(\Xi)\) is the event in the statement, then
\[
 \mathbb Q_m(E(\Xi))\le\mathbb Q_m(E^*)+\alpha_m.
\]
For a second independent displacement \(\Xi'\), the conditional
probability of being regular with jump error \(o(N)\) and having
\(\Delta(\Xi')\le q_m\) is at least \(1/2-\alpha_m\).
On \(E^*\), this implies the cycle event for \(\Xi'\) with
another \(q_m=o(N)\) exceptions.  Hence
\[
 (1/2-\alpha_m)\mathbb Q_m(E^*)
 \le\mathbb Q_m\{E'(\Xi'),\ \Xi'\text{ regular}\}.
\]
The preceding exponential estimate applies to the right side.
This proves the lemma.
\end{proof}

For comparison, the many-base result already in
\cite[Proposition~10.18]{P} uses fixed \(1<a<b\) with \(a-1\)
sufficiently small relative to
\((\vartheta-1)(b-a)\), and then a sufficiently large fixed
\(h_{\mathrm b}\).  If those bases lie in a \(C r_m(N)\)
interval and are separated by more than the largest high-range
cell width, the same all-but-\(o(N)\) cycle event has probability
tending to zero under the selected protocol.  We use this source
result in a narrow range and Lemma~\ref{lem:one-base-credit} in
the wide range.  The source's additional restriction on \(b-1\)
for its distance-length recursion is not a restriction of its
address or finite-view inputs.

\subsection{Circuits with controlled traffic and support}

We next convert an acyclic contour view into a flow of surrounding
circuits.  We state the geometric part of the finite-view input
separately from this new electrical conclusion.

\begin{lemma}[Retained data in an acyclic finite view]
\label{lem:finite-view-input}
Use the selected protocol, the packing bound
\eqref{eq:truncated-packing}, and the sparse all-seed cutoffs of
Proposition~\ref{prop:address-input}.  Fix \(1<a_0<a<b\).
Select and label the bases \(1,\ldots,h_{\mathrm b}\), and fix a
deterministic label \(\iota\in\{1,\ldots,h_{\mathrm b}\}\), before the
independent seed-generation mesh and address grids.  The common all-seed
trial array uses all distinguished bases and is the same for each fixed
choice of \(\iota\); the conclusions below concern base \(\iota\).
Use the finite-view construction with a uniform
seed in \([aN,bN]\), freeze its centers, and then use its
independent coarsening of log length \(L_m\to\infty\).  Its seed
separation parameter \(\mathsf H_m\to\infty\) is chosen so that
\(L_m=o(N\wedge\log\mathsf H_m)\).
Write \(X\) for the coarsened level.  Removing an event of
probability \(o(1)\), \(X\in[aN,bN]\), and its density on that
range is at most \(C/N\).  Retain the center graph and the acyclicity
indicators of its base components in the common trial-array extraction,
and follow the indicator for base \(\iota\); do not condition the view
law after this sampling.

After a further joint extraction, the following holds on that
acyclic event, with the fixed truncations and subsequent removal
of losses specified below.
Write \(r=r_m(X)\) and \(c_{m,i}\) for the center origins.
Two origins belong to the same scalar class of color \(c\) when
\[
 \frac{Z^c(c_{m,i})-Z^c(c_{m,j})}{\sqrt r}\longrightarrow0,
 \qquad
 \liminf_m\frac{\min_{[c_{m,i}\wedge c_{m,j},\,c_{m,i}\vee c_{m,j}]}
 Z^c-Z^c(c_{m,i})}{\sqrt r}\ge0.
\]
The combined graph is the union of these two colored class graphs.
\begin{enumerate}[label=(\roman*)]
\item There are finitely many center lines in the base's combined
component, with tight count, and distinct origins diverge in
\(r_m(X)\) units.  The base has a tight offset from its
line origin in \(r_m(X)\) units.  Consecutive lines of each scalar
class are joined by simultaneous actual scalar return links whose
endpoints approach their origins in these units.  For each such
colored center pair, the preselected witness graph contains the same
pair bit, with an actual seed witness whose endpoints have tight
coordinates in seed units.
\item A planar disk and an annular collar in an auxiliary sphere
agree with the original primal map on all their incidences and
complete stars needed for separation.  Their full time preimages
lie in retained intervals on the component's lines.  The inner disk contains the base walk
for a time radius \(\delta_m r_m(X)\), and the full preimage of
the closed outer disk has total length at most \(M_m r_m(X)\).
Here \(\delta_m^{-1}\), \(M_m\), and the number of lines are tight.
The maximum absolute coordinate of every retained interval endpoint
relative to its assigned line origin is tight in \(r_m(X)\) units.
For each fixed recipe truncation every retained slot and safety
band has duration between fixed positive constants in view units.
\item The annular collar admits a recipe in a countable family.
For a fixed recipe its input consists of retained off-center
block paths, forest shifts, durations, class anchors, and fixed
instructions.  Feasibility is measurable from that input.  Its
recorded port sets and confinement graphs use exact retained
incidences; every allowed path between its ports lies in the
prescribed outer tube.  The cyclic transverse tube pattern
forces any collection of successful recorded crossings to
contain a simple circuit in the original patch whose interior
contains the inner disk and whose closed disk lies in the outer disk.
The recipe has a positive-density ordinary sphere completion
of fixed positive duration after a fixed duration truncation.
\end{enumerate}
All slots and guards used by the annular recipe avoid the
distinguished center times.  Its observation family is a
countable family of fixed-recipe outputs chosen before the
randomized class anchors; the anchors can enter feasibility and
the later recipe selection.
\end{lemma}

\begin{proof}
The simultaneous saturation and all-seed construction in
\cite[Section~11]{P} uses precisely
\eqref{eq:truncated-packing}, the sparse cutoffs, and the
preselection event.  Its order is as follows.  First sample the
independent log-mesh phase, taking \(L_m=q_m\ell_R\) with integer
\(q_m\to\infty\) as in the source; changing the chosen length by
at most \(\ell_R\) preserves all displayed asymptotic conditions.
Run the truncated recipe at every
complete-cell endpoint.  Let \(P_N^{\rm seed}=\Theta(N/L_m)\)
count these cells and let \(\Lambda_m\) bound one truncated seed list.
Before the chronological grid is drawn,
round the seed depths as in the source, leaving them distinct,
and apply the terminal base-list augmentation at \(J_R\).
Coding then uses \(P_N\le P_N^{\rm seed}+1\) and
\(K_m=O(\Lambda_m+h_{\mathrm b})\); the two sparse ratios remain
vanishing because \(h_{\mathrm b}\) is fixed.  The independent
mesh phase is held as side information in this count.
Next draw the independent chronological
grid, form one forest from all the lists, and round and freeze
origins inside their coded terminal cells.  Then choose uniformly
among the original \(P_N^{\rm seed}\) complete cells, including
fallback cells, and
coarsen independently.  This law differs from the stated uniform
single-seed law by \(O(L_m/N)\) in total variation; the fallback
probability at that sampled cell vanishes along the sparse diagonal.
There is no conditioning that every cell succeeded.  Conditional
on the independent phases, the chosen seed-generation trial is
uniform among
\(P_N^{\rm seed}q_m=\Theta(N/\ell_R)\)
trials, including empty and
deleted trials.  For each of finitely many fixed residue classes
of sufficiently spaced generations, the retained enlarged raw
slots are disjoint after deletion.  The deletion probability tends
to zero by taking the discrete limit first and then the fixed
ancestor padding to infinity.  The address occupancy estimate
gives an expected code length \(O(N)\) under~\eqref{eq:selected-density}.
Entropy allocation in each residue class, followed by their finite
sum, gives bounded mean conditional entropy for this trial.
Its preliminary selected-slot comparison is
only a statement about continuum contour paths and finite height
differences.  It gives the stationary finite-line system without
asserting a law for discontinuous lattice outputs.
\cite[Lemma~11.12]{P} supplies the simultaneous actual return
links and whole-gap margins.  \cite[Lemma~11.24]{P} supplies the
identical original and auxiliary patch, including full stars and
full preimages.  The inner and outer time-preimage bounds in (ii)
are verified from this patch in the proof of
\cite[Proposition~11.30]{P}.
The patch construction first bounds all relevant partner coordinates
in a centered range \([-M,M]\) about their enrolled origins.
The retained rational cuts are chosen about those origins, and a finite
recipe truncation bounds the absolute values of their endpoint coordinates
in view units.
Finally \cite[Lemma~11.27]{P} constructs the countable
retained-data recipes and their open fillability.  The fixed-duration
ordinary completion is the separate construction immediately
preceding \cite[Proposition~11.30]{P}.
The source construction samples the class anchors after its
fixed-recipe observation vector.  All tight quantities can first
be bounded by fixed constants with an arbitrary fixed loss; its
finite data and sparse diagonals are then exhausted.
\end{proof}

Fix \(d_m,\ell_m\) with \(\ell_m\ge4d_m\), a tested chronological
window, an open planar confinement \(U\), and a depth \(D\).
A step occurrence of an original edge means either of its two
matched contour-step labels in the primal-tree and complementary-edge
dictionary.
A probability law \(\pi\) on simple original-edge circuits is
\emph{available at \(t\)} if every circuit in its support has
the following properties:
\begin{enumerate}[label=(\alph*)]
\item its open bounded interior contains the interpolated walk
on \([t-d_m,t+d_m]\);
\item its closed disk contains no uninterrupted run of
\(\ell_m\) consecutive walk steps in the tested window;
\item it is contained in \(U\), and each of its edges has a
step occurrence reached from \(t\) by at most \(D\) scalar links
or jumps of size at most \(\ell_m\).
\end{enumerate}
Here a simple circuit is an embedded Jordan circuit; a loop or
a two-edge circuit is allowed.
Its traffic is
\(\tau_\pi(e)=\pi\{C:e\in C\}\); circuits are simple, so each
edge is counted at most once.  At input level \(N\) we require
\begin{equation}
\label{eq:available-traffic}
 d_m=r_m(bN),\quad \ell_m=r_m(a_0N),\qquad
 \sum_e\tau_\pi(e)^2\le e^{\zeta N},
\end{equation}
with integer rounding and a fixed \(\zeta>0\).  The inequality
\(\ell_m\ge4d_m\) holds eventually.  A time is unavailable when
there is no such law.  The definition uses one fixed larger
window for the run test, and all bases and enlargements have
the prescribed spatial clearance there.

\begin{proposition}[Flow availability under selected laws]
\label{prop:flow-availability}
Fix \(\zeta>0\), a tested chronological window, and an open planar
confinement \(U\) with the preceding fixed-window and positive-clearance
cutoffs.  In the preceding definition of availability take \(D=D_m\)
and take \(d_m,\ell_m\) to be the integer-rounded values
in~\eqref{eq:available-traffic}.  Let \(D_m\to\infty\) be deterministic
and below the preselection
depth bound, with any finite-view cutoffs slowed below it.
Under~\eqref{eq:selected-units}, there is no sequence of selected
laws supported on \(h_{\mathrm b}\) unavailable times, pairwise
separated by \(\ell_m\), in an interval of length at most
\(C r_m(N)\), in either of these fixed parameter regimes:
\begin{enumerate}[label=(\roman*)]
\item \(1<a_0<a<b\), with \(a-1\) sufficiently small relative
to \((\vartheta-1)(b-a)\) and \(h_{\mathrm b}\) sufficiently
large as in \cite[Proposition~10.18]{P};
\item \(1<a_0<a\le2\), \(b=B\) from
Lemma~\ref{lem:one-base-credit}, and \(h_{\mathrm b}=1\).
\end{enumerate}
The same statement holds for every fixed multiple of the interval
length and every fixed positive confinement clearance.
\end{proposition}

\begin{proof}
Suppose such selected laws exist.  Mark and label the alleged unavailable
bases before the independent seed mesh, address grid, and view sampling,
and fix an arbitrary
deterministic label \(\iota\) at that stage.  Let \(\mathsf S_m\) denote
the selected sampled-trial law on the common array and its successive
extensions by the stated construction kernels for base \(\iota\).
Let \(\mathsf A_m^{\mathrm{acyc}}\) be the retained indicator from
Lemma~\ref{lem:finite-view-input}, belonging to the common trial marginal,
that base \(\iota\)'s combined center component is acyclic, and write
\(\mathsf A^{\mathrm{acyc}}\) for its extracted limit.  The auxiliary
extensions may depend on \(\iota\); no independence between extensions
for different labels is used.

Use Lemma~\ref{lem:finite-view-input} and restrict first to
finitely many recipe indices and fixed duration, line-count, and
margin bounds, losing an arbitrary fixed probability.  For each
fixed recipe, before its class anchors are drawn, record
\[
 Y_{m,s}=\EL(\mathcal P_{m,s}),\qquad 1\le s\le q,
\]
in raw unit-edge units, where \(\mathcal P_{m,s}\) is its full
family of paths between the recorded vertex ports in the
recorded confinement graph.  Use \(+\infty\) for an empty family
and compactify these coordinates.  This is a function of the
same exact incidence record and fixed intrinsic instruction,
as required by Proposition~\ref{prop:readout-extension}.

For each fixed recipe and fixed truncation, consider the compatible
anchored selected subprobabilities constructed below, retaining
\(\mathsf A_m^{\mathrm{acyc}}\) as an extra coordinate.  Let \(\nu\)
be any weak limit of their pushforwards to normalized view coordinates,
after the nuisance trial, unit, and parity labels are projected away
at the step specified below.  Write \(c_0\) for the limiting intrinsic
pre-anchor input, \(a\) for its class anchors, and \(y=(y_s)_{s=1}^q\)
for the extended-real electrical output in \([0,\infty]^q\), retained
through compactification.  Let \(F(c_0,a)\) be the
fixed recipe's anchored feasibility event in the retained limiting
view of Lemma~\ref{lem:finite-view-input}; its strict margins give
the eventual lattice crossing interpretation.  Our target is
\begin{equation}
\label{eq:selected-route-null}
 \nu\{\mathsf A^{\mathrm{acyc}}=1,\ F(c_0,a),\
       \max_{1\le s\le q}y_s=+\infty\}=0.
\end{equation}
This is a null statement for an unnormalized subprobability; no
acyclic, compatibility, or successful-recipe probability is divided out.

In an ordinary completion of the fixed duration, feasibility has
a countable family of rational route witnesses.  Work first with
one such witness: choose compact endpoint balls inside its open
ports and a compact connecting route inside its recorded open
confinement.  Finitely many ordinary charts off the auxiliary
marks cover it.  Theorem~\ref{thm:network-bounds} gives a path
law with finite squared traffic there.  Full preimages and the
strict inner/outer margins put its endpoints in the recorded
ports and every used edge in the recorded graph for all large
indices.  The countable intersection covers every feasible
witness.  Therefore the
limit of \(Y_{m,s}\) is finite on feasibility in every jointly
extracted ordinary completion.  This uses only the fixed-test
electrical conclusion, for arbitrary diverging units and a
fixed duration; neither its threshold nor its route constant
has to be uniform over recipes.

We spell out the transfer to the selected view.  Let \(c_{0,m}\)
be the retained input before the class anchors are drawn, including
the block paths, forest shifts, durations, fixed instructions, and
the trial, unit, and parity parameters \(\theta_m\).  Let \(Y_m\)
be the fixed electrical subvector above.  In each residue trial
the code and recipe fix the requested cuts, durations, and guard/noise
search anchors before fresh letters; auxiliary choices have total
conditional weight at most one.  These search anchors precede the
class anchors considered below.  The shifts form the prescribed
forests, and fixed truncation supplies positive safety bands and
compact fitting ranges.  Let \(G_m^0\) be genuine raw-slot
compatibility for this pre-anchor instruction.  Use the pre-anchor
marginal of \(\mathsf S_m\) in the following comparison, without
conditioning on acyclicity or recipe success.

Write \(\mathsf R_m^0(dc_0,dy)\) for the independent raw-block
reference law of this pre-anchor input and output, integrated
against the same retained trial and parameter marginal as
\(\mathsf S_m\), and set
\(\mathsf S_m^0(B)=\mathsf S_m\{G_m^0,(c_{0,m},Y_m)\in B\}\).
The finite domination part of Lemma~\ref{lem:record-null-transfer},
followed by the bounded all-seed trial entropy, gives
\begin{equation}
\label{eq:pre-anchor-tac}
 \mathsf R_m^0(B_m)\longrightarrow0
 \quad\Longrightarrow\quad
 \mathsf S_m^0(B_m)\longrightarrow0.
\end{equation}
This is tight absolute continuity of the compatible
\emph{subprobability} law; no division by the compatibility mass
occurs.  The comparison sees the lattice output \(Y_m\), whereas
the preliminary comparison in Lemma~\ref{lem:finite-view-input}
only concerned continuum contours.

After the entire pre-anchor selected sample, including its full
fixed observation vector, draw the finite class-anchor vector
\(a_m\).  Its translation may depend on \(Y_m\).  First restrict
to finitely many anchor counts and compact height ranges, keeping
these restrictions unnormalized.  Its conditional density given
that entire sample is bounded, as in \cite[Lemma~11.25]{P}.
Choose a positive smooth reference anchor density, with its scaled
feasible-lattice version \(\varrho_{m,\theta_m}(da)\).  The lattice
law is expressed relative to covolume-normalized counting measure;
on the compact truncation its density is uniformly bounded below.
On those ranges the restriction of the actual conditional anchor
law is therefore bounded above by \(C\varrho_{m,\theta_m}\).
The same bound given \((c_{0,m},Y_m)\) follows by conditional
expectation.  The reference anchors are independent
of \(Y_m\) conditional on the recorded \(\theta_m\).

Here is the sectional argument, applied without normalizing
\(\mathsf S_m^0\).  Let \(\mathsf S_m^a\) be the joint selected
subprobability after adjoining \(a_m\), restricted to \(G_m^0\)
and the chosen anchor-count and height ranges,
and put
\[
 \mathsf R_m^a(dc_0,dy,da)
   =\mathsf R_m^0(dc_0,dy)\varrho_{m,\theta(c_0)}(da).
\]
For an event \(B_m\) with \(u_m=\mathsf R_m^a(B_m)\to0\), let
\(p_m(c_0,y)=\varrho_{m,\theta(c_0)}
 \{a:(c_0,y,a)\in B_m\}\).
Markov's inequality gives
\(\mathsf R_m^0\{p_m>\sqrt{u_m}\}\le\sqrt{u_m}\), while the
conditional anchor bound gives
\[
 \mathsf S_m^a(B_m)
 \le \mathsf S_m^0\{p_m>\sqrt{u_m}\}+C\sqrt{u_m}
 \longrightarrow0
\]
by~\eqref{eq:pre-anchor-tac}.  This is the proof of
\cite[Lemma~11.25]{P}, including its scaled-lattice form.
It allows the actual anchors to depend on the output.  Apply it
separately to each fixed finite recipe vector; the countable null
intersection and the preceding finite recipe truncation handle
the later recipe selection.

Before taking the kernel limit, project away the trial index and
raw unit/parity labels after expressing the retained paths, shifts,
durations, and anchors in normalized view units.  Tight absolute
continuity survives this pushforward.  The fiberwise anchor reference laws
are lattice discretizations of one fixed smooth law \(\varrho\);
on each normalized compact they converge to \(\varrho\) uniformly
over the allowed unit and parity parameters as the minimum unit
diverges.  The normalized reference factorization of
\cite[Theorem~7.10]{P} and Proposition~\ref{prop:readout-extension}
therefore gives the independent-anchor limit law
\(\mu_0(dc_0)K_{c_0}(dy)\varrho(da)\).  The positive ordinary
filling argument, applied to augmented input \((c_0,a)\), gives
the retained subprobability law
\[
 f(c_0,a)\,\mu_0(dc_0)\varrho(da)K_{c_0}(dy),
 \qquad f(c_0,a)>0\quad\text{on }F.
\]
As in~\eqref{eq:fillable-density}, \(f\) integrates the strict
omitted filling and its input-dependent density weight, so it has
no \(y\) argument.  The ordinary finiteness proved above yields
\[
 0=\int f(c_0,a)\mathbf1_{F(c_0,a)}
 K_{c_0}\{\max_s y_s=\infty\}\,
 \mu_0(dc_0)\varrho(da).
\]
Positivity of \(f\) makes the anchored feasibility/infinite-cost
event null for the independent-anchor reference product.
The post-anchor tight absolute continuity just proved transfers
this null event to the selected anchored limit on compatibility.
Every successful selected recipe has \(G_m^0\) by its construction
from the actual compatible raw instructions.
The projection of \(\nu\) restricted to
\(\{\mathsf A^{\mathrm{acyc}}=1\}\) is a submeasure bounded by this
compatible anchored marginal, so~\eqref{eq:selected-route-null}
follows without a kernel or independence assertion for the indicator.
Only the independent-anchor reference and ordinary filling laws
use the displayed output kernel \(K_{c_0}\); the selected anchored
law is used through tight absolute continuity.

All comparisons through the anchor step are finite lattice
comparisons at fixed recipe, guard, duration, shift, anchor-count,
and height truncations.  Their effective slot lengths diverge
because \(r_m(bN)\to\infty\).  Apply them in each joint extraction
with \(r_m(X)\) and its recorded unit parameters; the normalized
reference construction and the scaled-lattice anchor argument
apply to these varying laws.  The ordinary finiteness assertion
holds for every such extraction, so the argument also covers
mixtures over \(X\), as in \cite[Proposition~7.18]{P}.
The anchor restrictions are exhaustible under the selected law:
on a fixed line-count and centered-oscillation truncation, the
completion assigns class heights
\(H_A=(4K_0+10)s(A)+u_A\), with \(1\le s(A)\le k\) and
\(\abs{u_A}<1\), where \(K_0\) bounds the retained centered
oscillations.  Bounding \(k\) and \(K_0\) bounds these normalized
heights; bounded density alone is not used for their tightness.
Fix the finite recipe list and the other stated truncations.  By the
strict certificate limits in \cite[Lemma~11.27]{P},
\eqref{eq:selected-route-null} gives vanishing unnormalized selected
mass for the finite-threshold cost failure of a successful retained
recipe, first as \(m\to\infty\) and then as the cost threshold tends
to infinity.  Only afterward exhaust the finite recipe and truncation
losses.  No slot meets the center, and feasibility has no omitted
filling as an argument.  This is the restricted tightness of
\(\sum_sY_{m,s}\) used below.

By Lemma~\ref{lem:finite-el}, choose a near-minimizing path law for each \(\mathcal P_{m,s}\)
with squared traffic at most \(Y_{m,s}+1\).  Sample those paths
and select a simple surrounding circuit from their union.
Every sample works by the recipe's transverse pattern.  If its
traffic is \(\tau\) and the path traffics are \(\tau_s\), then
\[
 \tau(e)\le\sum_{s=1}^q\tau_s(e),\qquad
 \sum_e\tau(e)^2\le q\sum_{s=1}^q\sum_e\tau_s(e)^2.
\]
On each fixed truncation, its unnormalized upper tail is controlled
by the restricted cost estimate.  Since \(e^{\zeta N}\to\infty\),
the corresponding traffic-failure mass tends to zero.

The inner and outer preimage bounds in
Lemma~\ref{lem:finite-view-input}(ii) give the two disk
conditions.  The event \(bN-X\le\sqrt N\) has probability
\(O(N^{-1/2})\).  Outside it,
\(\delta_m r_m(X)/r_m(bN)\to\infty\) after the fixed tightness
truncation.  Also
\[
 \frac{M_mr_m(X)}{r_m(a_0N)}
 \le M_m e^{-(a-a_0)N}\longrightarrow0.
\]
Every uninterrupted visit in a closed circuit disk, including
a run on its edges, belongs to the full outer preimage, so it
has length less than \(\ell_m\).  These conclusions apply to
every circuit in the constructed law.

They also have the required support reach from the \emph{given}
base.  If its combined component has \(K\) lines, a simple path
of consecutive scalar-class connections from the base line to
any retained line uses at most \(K-1\) actual return links.
Interleaving jumps on those lines and a final jump to a retained
edge occurrence uses at most \(2K+O(1)\) steps, each jump bounded
by a tight multiple of \(r_m(X)\).  The same finite return
witnesses work for every occurrence in every retained interval.
All path support is in the protected patch, whose complete time
preimages are on those intervals.  Since
\(r_m(X)/r_m(a_0N)\le e^{-(a-a_0)N}\), the jumps are eventually
at most \(\ell_m\), and tight \(K\) is eventually below the
diverging \(D_m\).  The master event and its complete-star mesh
bound put the entire supported edges in the required confinement.
Combining these vanishing disk and support errors with the restricted
cost estimate and then exhausting the fixed losses gives
\[
 \mathsf S_m\!\left(
 \{\mathsf A_m^{\mathrm{acyc}}=1\}\cap
 \{\text{base }\iota\text{ is unavailable}\}
 \right)\longrightarrow0.
\]
Here unavailable means failure of a law satisfying (a)--(c) and
\eqref{eq:available-traffic}.  Under the selected laws assumed for
contradiction, base \(\iota\) is unavailable with probability one, hence
\(\mathsf S_m\{\mathsf A_m^{\mathrm{acyc}}=1\}\to0\).

It remains to use the address charge.  The bases in (i) are
\(\ell_m\)-separated, whereas high-range cell widths are
\(O(r_m(aN))=o(\ell_m)\); the same remains true along the
permitted sparse enclosure diagonal.  The later return links
certify colored center-pair bits in the saturated graph.  For each
such bit the all-seed code retained an earlier chosen actual seed
witness and jumps associating its endpoints to the corresponding
original pre-grid seed centers.  The association jumps cost \(o(N)\)
split generations.  A base has a tight seed offset from its
original pre-grid seed center.  On a fixed offset truncation, the independent-grid
probability that they are separated after coarsening by \(U\)
is \(O(\min\{1,e^{-U}\})\), whose average for
\(U\sim\operatorname{Unif}[0,L_m]\) is \(O(L_m^{-1})\).
Thus only \(o(N)\) further trial generations fail, in probability.
Rounding an origin preserves its seed center's terminal cell and all its ancestors.
Outside these exceptions the centers, witness endpoints, and
bases have the corresponding projected cells.
Thus a rank cycle of a center component is a rank cycle in the
projected base component.  Adding further connected vertices
cannot destroy it, since a union spanning tree gains at least
one total monochrome rank per added vertex.  Apply this argument
separately to each fixed label \(\iota\) on the common trial array.
Let \(E_{m,\iota}\) count generations
\(k\ell_R\in[aN,bN]\) at which its projected base component lacks a
rank cycle.  After including the conversion exceptions above, averaging
the vanishing acyclic probability over the uniform seed-generation
trials and applying Markov's inequality gives
\(E_{m,\iota}/N\to0\) in probability under the common selected array
law.  Since \(h_{\mathrm b}\) is fixed,
\(N^{-1}\sum_{\iota=1}^{h_{\mathrm b}}E_{m,\iota}\to0\) in that same
probability.  The number of generations exceptional for at least one
label, divided by \(N\), is bounded by this normalized sum.  Outside
that union every projected base component has a rank cycle.
Choose a deterministic
\(\eta_m\downarrow0\) so that the exceptional count is at most
\(\eta_mN\) with probability tending to one.  This is the cost-independent
witness conversion in the proof of the circuit-sparsity result
in \cite[Section~12]{P}.

In regime (i) this contradicts
\cite[Proposition~10.18]{P}.  In regime (ii) it contradicts
Lemma~\ref{lem:one-base-credit}.  The latter uses only fixed
\(B\) and \(r_m(BN)\to\infty\); the distance normalization and
the narrow length-recursion restriction in~\cite{P} were not
used.  Removing the arbitrary finite truncation losses proves
the proposition.
\end{proof}

\subsection{Uniform levels and circuit covers}

We pass from the selected-law contradictions to one event for
many levels.  Let the tested levels lie in a finite union of
fixed-spaced lists and their fixed positive multiples, all
starting at \(N_0(m)\to\infty\).  Assume the minimum units in
each used view diverge uniformly.  Choose the master event
before selecting any level, with
\(\sigma_m\) a fixed multiple of \(e^{-N_0(m)}\).

If a union of failures has probability at least \(\delta>0\),
take the master loss smaller than \(\delta/4\).  Since the sum
of \((N+1)^{-2}\) over these lists is bounded, some deterministic
level \(N\) has failure together with that same master event
of probability at least \(c\delta(N+1)^{-2}\).  Conditioning
on this event gives raw-word density at most its reciprocal:
indeed the conditional event probability given the word is at
most one.  Its logarithm is \(O_\delta(\log(N+1))=o(N)\).
Choose the failing bases or truncated witnesses afterwards and
the independent grids last.  Proposition~\ref{prop:flow-availability}
or Lemma~\ref{lem:reachable-packing} contradicts this selection.
This is the weighted-level argument of
\cite[Lemma~12.5]{P}, now applied to fixed-spaced lists as well.

For packing, do this first for every fixed reach depth.  A
deterministic diagonal then gives a depth tending to infinity
slowly enough that all these union probabilities still vanish,
and still below the master bound.  Use a fixed fraction of it
for circuit support, leaving room for two access chains and a
bounded number of extra steps.  Proposition~\ref{prop:flow-availability}
allows this slower diverging depth.  A finite list of scale
multiples is handled by the same diagonal.  Exhausting the
fixed master and clearance losses, we henceforth work on events
of probability \(1-o(1)\) where these conclusions hold at
every used level.  This uniformity comes from the address
argument; it is not an assertion about the thresholds in
Theorem~\ref{thm:network-bounds}.

Here is the resulting overlap bound.  At a primary level \(N\)
put \(T=r_m(N)\).  Consider disjoint blocks of lengths between
fixed multiples of \(T\), and a fixed-multiple enlargement
\(\widehat I\) of each.  For each available test which will be
used from a time in \(\widehat I\), choose one witnessing flow
independently of \(f\).  Let \(W_I\) contain all edges in those
flow supports and all chronological walk edges in \(\widehat I\).
The recursive tests below have input \(y\ge N\), so their
support jumps are at most \(r_m(a_0y)\le r_m(a_0N)\); a wide
test started at \(N\) has the same bound.  The master event
and the fixed clearance put these sets in \(G_m^+\).

\begin{lemma}[Overlap of the energy sets]
\label{lem:energy-overlap}
At every such level with uniformly diverging view units, the
common event above can be chosen so that, for every family of
disjoint blocks and every vertex function,
\begin{equation}
\label{eq:energy-overlap}
 \sum_I\mathcal E_m(f;W_I)
 \le C e^{N/40}\mathcal E_m(f;G_m^+).
\end{equation}
At levels where only chronological edges are used, the same
statement holds with a constant in place of \(Ce^{N/40}\).
\end{lemma}

\begin{proof}
Fix an edge \(e\) belonging to several \(W_I\).  From those
blocks retain a sublist whose enlargements are separated by a
fixed multiple of \(T\), losing only a fixed factor.  For each
retained block choose the responsible time in \(\widehat I\);
it reaches a step occurrence of \(e\) with the chosen depth and
jump bound, or is itself such an occurrence for a chronological
edge.

By the dictionary of incidence occurrences in
\cite[Proposition~3.1 and Lemma~3.2]{C}, each original edge has
the two contour steps of one matched burger/order pair:
\(L\)-pairs represent primal-tree edges, while the quadrangles
of \(R\)-pairs carry the complementary primal edges.  The
convention includes loops, bridges, and repeated incidences.  If the
up-step ends at \(u\) and its matching down-step ends at \(v\),
their scalar baselines satisfy
\[
 Z(u-1)=Z(v)\le \min_{u-1\le s\le v}Z(s).
\]
Thus \((u-1,v)\) is an actual scalar link; either step occurrence
is within one time step of the corresponding endpoint.
One link and at most two unit jumps join its two occurrences,
including for loops and parallel edges.  The units
\(r_m(a_0N)\) diverge whenever a circuit is used, so these unit
jumps are allowed.  Joining through \(e\) therefore reaches
every responsible time from the first within twice the
support depth plus a fixed number of steps.  The times are
\(cT\)-separated.  The simultaneous version of
Lemma~\ref{lem:reachable-packing}, with \(\zeta=1/40\), bounds
their number by \(e^{N/40}\).  This bounds the multiplicity
of each edge in the left side.

For chronological-only sets, the enlarged intervals have bounded
time overlap, and the primal tree tour traverses a tree edge at
most twice.  Their edge multiplicity is bounded.  Summing the
edgewise bounds proves the lemma.
\end{proof}

We use the following deterministic fact about a nearest-neighbor
walk \(v\) in a planar graph, interpolated along its edges.
For a circuit \(C\), put
\(\operatorname{var}_f(C)=\sum_{e\in C}\abs{df(e)}\).

\begin{lemma}[Covering and joining a chronological segment]
\label{lem:segment-cover}
Let \(d\ge1\), \(\ell\ge4d\), and \(h_{\mathrm b}\ge1\) be
integers.  An integer \(t\) is called good if it has a simple
circuit in \(U\) with \(\operatorname{var}_f(C)\le L\), whose
open interior contains \(v([t-d,t+d])\), and whose closed disk
contains no run of \(\ell\) steps in the tested window.
There is a constant \(K_{h_{\mathrm b}}\), depending only on
\(h_{\mathrm b}\), with the following property.  Suppose an interval
\(I\) of at most \(T\) steps has no
\(h_{\mathrm b}\) pairwise \(\ell\)-separated bad times, and
the walk on a \(K_{h_{\mathrm b}}\ell\) enlargement stays in
\(U\), with that enlargement inside the tested window.  Then:
\begin{enumerate}[label=(\roman*)]
\item At most \(K_{h_{\mathrm b}}(1+T/d)\) good circuits cover
the walk on \(I\) by their open interiors except for at most
\(K_{h_{\mathrm b}}\) chronological intervals, each of length
at most \(K_{h_{\mathrm b}}\ell\).
\item There are graph paths in \(U\) with total \(\abs{df}\)-cost
at most \(K_{h_{\mathrm b}}(1+T/d)L\) which, together with at
most \(K_{h_{\mathrm b}}\) unresolved endpoint pairs separated
by at most \(K_{h_{\mathrm b}}\ell\), join the endpoints of
the walk on \(I\).  All unresolved times lie in that
enlargement.
\end{enumerate}
\end{lemma}

\begin{proof}
A maximal \(\ell\)-separated set of bad times has at most
\(h_{\mathrm b}-1\) elements.  Their \(\ell\)-neighborhoods
cover the bad times and, after merging, have total length
\(O(h_{\mathrm b}\ell)\).  The complement is a union of at
most \(h_{\mathrm b}\) good integer runs.  In each run choose
its endpoints and good times spaced at most \(d\).  Their
number is at most \(2+T/d\) per total run length, and their
open interiors cover the complete interpolated run.  Enlarging
the deleted intervals by a bounded number of steps handles
the endpoints and proves (i).

For (ii), in a good run retain the chosen disks maximal for
inclusion.  Consecutive original interiors overlap, so the
maximal interiors have connected intersection graph.  Two
overlapping incomparable Jordan disks have intersecting
boundaries.  Their boundary union is consequently a connected
subgraph with total \(\abs{df}\)-cost at most the sum of the
chosen circuit costs.  The run's first vertex reaches the
boundary of a disk containing it within \(\ell+1\) forward
steps, and its last vertex does so within \(\ell+1\) backward
steps, by the no-run condition.  A hit is at a graph vertex,
since the walk and circuit use the same embedded graph.
Join the two hits in the boundary union.  Concatenating these
connections across the good runs leaves only the deleted
intervals and at most two such exit lengths at each gap.
Their number and lengths have the asserted bounds.  If there
is no good run, use the one unresolved endpoint pair.  This
is the additive-cost proof of \cite[Lemma~12.2]{P}; it applies
to the nonnegative edge costs \(\abs{df}\).
\end{proof}

Choose the scale constants now.  Take \(B\) from
Lemma~\ref{lem:one-base-credit}, choose
\[
 0<u<10^{-5}/B,\qquad b=1+u,
\]
and then choose
\[
 1<a_0'<a_0<a<b
\]
with \(a-1\) small enough for regime (i) of
Proposition~\ref{prop:flow-availability}.  Fix the corresponding
\(h_{\mathrm b}\).  Use \(\zeta=u\) both in this narrow test
and in the wide test with upper endpoint \(B\).  These choices
are fixed.  Each narrow available law at input \(y\), supported
in \(W_I\), supplies a circuit with
\begin{equation}
\label{eq:one-circuit-bill}
 \operatorname{var}_f(C)
 \le e^{uy/2}\sqrt{\mathcal E_m(f;W_I)}.
\end{equation}
Indeed the expected variation is the inner product of traffic
with \((\abs{df(e)})_e\), so Cauchy--Schwarz gives this bound
for at least one circuit.  This circuit may depend on \(f\);
the support set \(W_I\) does not.

Consider first
\begin{equation}
\label{eq:coarse-range}
 N\le (4B)^{-1}\log m.
\end{equation}
Apply Lemma~\ref{lem:segment-cover}(i) at input \(y=N\) to a
block \(I\) of \(O(r_m(N))\) steps.  Its circuits cover all
but a bounded number of intervals of length
\(O_{h_{\mathrm b}}(r_m(a_0y))\).  For large \(y\) each fits
in an interval of length \(r_m(a_0'y)\).  Apply the same
construction there with new input \(a_0'y\).  Continue until
the new input would be at least \(BN\).  This takes a fixed
number of stages, since \(a_0'>1\) and \(B\) are fixed.
At the last stage each remaining interval has length
\(o(r_m(BN))\).  The wide, one-base test started at \(N\)
is available at every time, so constantly many of its
circuits cover each remaining interval.

At a stage with input \(y\), the number of circuits per current
interval is \(O_{h_{\mathrm b}}(1+e^{uy})\), and its circuit
bill is bounded by~\eqref{eq:one-circuit-bill}.  The number
of current intervals grows by only a fixed factor at each
of a fixed number of stages.  Since every narrow input is
less than \(BN\), the total bill is at most
\[
 C_{B,u,h_{\mathrm b}}e^{5uBN}
       \sqrt{\mathcal E_m(f;W_I)}
 \le e^{N/10}\sqrt{\mathcal E_m(f;W_I)}
\]
for all sufficiently large \(N\).  The final wide circuits
have bill \(e^{uN/2}\sqrt{\mathcal E_m(f;W_I)}\) and fit the
same bound.  Every disk used has a run cap at most
\[
 r_m(a_0N)=e^{-(a_0-1)N}r_m(N)=o(r_m(N)).
\]

Discard any disk whose open interior misses the walk on \(I\),
and retain the disks maximal for inclusion.  An open cover of
a connected walk has connected intersection graph.  The
Jordan argument in Lemma~\ref{lem:segment-cover} therefore
makes the union \(S_I(f)\) of these maximal boundaries a
connected subgraph.  It follows that
\begin{equation}
\label{eq:boundary-oscillation}
 \osc(f;S_I(f))
 \le e^{N/10}\sqrt{\mathcal E_m(f;W_I)}.
\end{equation}
This is a cover of the whole walk by open interiors, with the
no-long-run property for every closed disk.  It has not paid
to join an individual interior vertex to a boundary.

For the remaining levels, \(N>(4B)^{-1}\log m\), we instead
obtain
\begin{equation}
\label{eq:point-oscillation}
 \osc_{i\in I}f(v(i))
 \le e^{N/10}\sqrt{\mathcal E_m(f;W_I)}.
\end{equation}
If \(N>(1-4u)\log m\), a chronological path in the block has
at most \(O(1+m^{4u})\) edges, with bounded multiplicity.
Cauchy--Schwarz gives a bill
\(O(m^{2u})\sqrt{\mathcal E_m(f;W_I)}\).
Otherwise apply Lemma~\ref{lem:segment-cover}(ii) at input
\(y=N\) to any two cuts to be joined.  Its bill is
\[
 C_{h_{\mathrm b}}(1+e^{uy})e^{uy/2}
       \sqrt{\mathcal E_m(f;W_I)},
\]
and only a bounded number of pairs of separation
\(O_{h_{\mathrm b}}(r_m(a_0y))\) remain.  Repeat at input
\(a_0'y\) until it exceeds \((1-4u)\log m\), and then use
chronological paths on the remaining pairs.  The number of
stages is bounded in terms of \(B,u,h_{\mathrm b}\), since
the ratio of the final input to the first is at most \(4B\).
The total circuit bill is at most
\(C m^{3u/2}\sqrt{\mathcal E_m(f;W_I)}\), and the terminal
bill at most \(C m^{2u}\sqrt{\mathcal E_m(f;W_I)}\).
Both are eventually bounded by the right side of
\eqref{eq:point-oscillation}, because
\(e^{N/10}\ge m^{1/(40B)}\) and \(2u<1/(40B)\).
All unresolved pairs stay in a fixed-multiple enlargement
of the original block.

These uses meet the diverging-unit condition for the common
level event.  In~\eqref{eq:coarse-range}, the finest narrow
level is less than \(2BN\) and the wide level is \(BN\), so
the minimum units are at least \(m^{1/2}\).  In the fine
recursion \(y\le(1-4u)\log m\), the finest narrow level is
\((1+u)y\); its minimum units are at least
\(m^{3u+4u^2}\).  The levels used are a fixed finite list
of multiples of the primary levels.  At still finer levels
only the chronological bound is used.  Thus
\eqref{eq:energy-overlap}--\eqref{eq:point-oscillation}
hold together for all levels starting at any deterministic
\(N_0(m)\to\infty\), with probability tending to one after
the fixed cutoffs are exhausted.

\subsection{The dyadic estimate and its threshold}

Let \(\mathcal K_m^{\mathrm{act}}\) denote a fixed activation cutoff
for a bilateral protected interval:
its cut length divided by \(m\) lies between fixed positive
constants, its labels lie in a fixed chronological window,
and its limiting image has a fixed positive margin inside
the intermediate neighborhood recorded by \(G_m^+\).
The preceding probability statements are uniform over all
interval starts and roundings on this cutoff.  Indeed their
unavailable-time events range over every interval in that
window and their packing bounds range over arbitrary lists.
The master event contains every fixed-multiple enlargement.

At dyadic level \(j\), use primary input \(N=j\log2\).
Nonterminal block lengths are between fixed multiples of
\(r_m(N)\), with bounded integer rounding.  For a block in
the regime~\eqref{eq:coarse-range}, choose \(F_I\) to be
\(f\) at any vertex of \(S_I(f)\).  For a block in the
pointwise regime choose its value at any occurrence in \(I\).
For a child \(I'\) of \(I\),
\begin{equation}
\label{eq:parent-child}
 \abs{F_{I'}-F_I}
 \le C e^{N/10}
 \left(\sqrt{\mathcal E_m(f;W_{I'})}
       +\sqrt{\mathcal E_m(f;W_I)}\right).
\end{equation}

Here is the needed comparison when boundaries are used.
Suppose both blocks have interior covers but their maximal boundary
unions are disjoint.  Fix one point of the child walk and choose
from each cover a disk whose open interior contains that point.
With disjoint boundary unions
the two disks would be nested.  The connected boundary union
of the inner family, avoiding the outer boundary, lies
entirely in the containing disk.  Every bounded Jordan disk
with boundary there also lies there.  Hence that family's
entire covered walk lies in one disk of the other family.
In either nesting direction this contains the child walk,
whose length is a fixed multiple of \(r_m(N)\), contrary
to the disk's \(o(r_m(N))\) run cap.  If the parent has a
boundary cover and the child is in the pointwise regime, a
child walk missing the parent boundaries would similarly
stay in a single parent disk, again impossible.  A boundary
hit includes a common graph vertex.  Comparing at that
vertex by~\eqref{eq:boundary-oscillation} and
\eqref{eq:point-oscillation} proves
\eqref{eq:parent-child}.  If both blocks are pointwise,
compare at any common child occurrence.  The cover regime
has diverging block lengths, so rounding cannot affect the
run-cap comparison.

Let \(g_j(i)=F_I\) on each level-\(j\) block.  Give each cut
mass \(1/m\).  A child block has mass \(O(e^{-N})\).
Squaring~\eqref{eq:parent-child}, summing its child and parent
energies, and using~\eqref{eq:energy-overlap} at the two
levels gives
\begin{equation}
\label{eq:successive-blocks}
 \norm{g_{j+1}-g_j}_{L^2(I_m,m^{-1})}^2
 \le C e^{(-1+1/5+1/40)N}\mathcal E_m(f;G_m^+)
 =C e^{-31N/40}\mathcal E_m(f;G_m^+).
\end{equation}
The last singleton function is exactly \(i\mapsto f(v(i))\).
The sum of the square roots of the right sides of
\eqref{eq:successive-blocks} is geometric in \(j\).
Minkowski's inequality therefore gives
\[
 \norm{f(v(\cdot))-g_j}_{L^2(I_m,m^{-1})}^2
 \le C'2^{-31j/40}\mathcal E_m(f;G_m^+).
\]
Orthogonal projection minimizes this error, so the same
bound holds with \(P_j^{I_m}f\) in place of \(g_j\).
Since \(31/40>p=7/10\), the constant \(C'\) can be absorbed
by increasing the starting level.  The auxiliary circuit
choices depended on \(f\), but \(W_I\) and their overlap
event did not, and the final projection is linear.  The
conclusion is simultaneous in all \(f\).

In particular, for every deterministic \(j_0(m)\to\infty\),
\begin{equation}
\label{eq:bilateral-all-level-failure}
 \PP\!\left(
 \mathcal K_m^{\mathrm{act}}\cap
 \left\{
 \begin{gathered}
 \exists j\ge j_0(m),\ \exists f:\\
 \frac1m\sum_{i\in I_m}
 \abs{f(v(i))-(P_j^{I_m}f)(i)}^2
 >2^{-pj}\mathcal E_m(f;G_m^+)
 \end{gathered}
 \right\}\right)\longrightarrow0.
\end{equation}
This follows from the common level event above;
if \(j_0\) is beyond singleton scale the assertion is
immediate.  Letting the arbitrary master losses vanish and
exhausting the fixed window, duration, and clearance cutoffs
gives the assertion locally on activation.

Define \(\mathfrak j_m\) to be one past the largest level at which
the failure inequality in~\eqref{eq:bilateral-all-level-failure} holds, with value
zero if there is none.  It is finite because the projections
are eventually the identity.  The sequence
\(\mathfrak j_m\) is tight on each fixed activation cutoff.
Indeed, if for some \(\delta>0\) it were not tight, one
could choose \(m_k\to\infty\) and \(J_k\to\infty\) with
\(\PP(\mathcal K_{m_k}^{\mathrm{act}}\cap
\{\mathfrak j_{m_k}>J_k\})>\delta\).
Choose a deterministic schedule \(j_0(m)\to\infty\) slowly
enough that \(j_0(m_k)\le J_k\).  This contradicts
\eqref{eq:bilateral-all-level-failure}.  Thus the threshold
is tight \emph{before} any sphere transfer.  Separate
estimates at each fixed \(j\) would not have implied this
conclusion.

\subsection{Transfer of the threshold to sphere charts}

We finish the proof of Theorem~\ref{thm:block-approximation}.
For one fixed protected instruction, compute
\(\mathfrak j_m\) from the exact graph \(G_m^+\), its
identified interval vertices, and its integer block partitions.
For a given \(j\), failure for some vertex function is a
finite-dimensional quadratic comparison: after fixing the
vertex labels it asks whether a specified symmetric matrix
is positive semidefinite.  Hence \(\mathfrak j_m\), compactified
in \([0,\infty]\), is a measurable numerical readout of that
finite record.  It uses relative cut labels and the common
deterministic unit \(m\), with no absolute tour origin,
omitted-gap duration, or unrecorded edge.

Choose its energy graph using an inner time-cell cover which
contains every star over a compact neighborhood of the
interval image, and whose closure is still in the larger
protected chart.  The paid \(W_I\) lie in that inner
neighborhood on the master events.  Exact reconstruction
with the intermediate margins identifies every tested
vertex and edge.  Adding other correctly reconstructed edges
only enlarges the right side.  Such choices form a cofinal
countable family by full-preimage compactness, and activation
depends only on the limiting local data and relative order.
The observation can have an arbitrary default off activation.

In every joint bilateral extraction the compactified limit of
\(\mathfrak j_m\) is finite on activation.  To see this
directly when the activation itself is random, suppose a
positive part of an activated event had limit \(+\infty\).
Restrict it to finitely many positive-clearance, duration,
and chronological cutoffs with still positive probability.
The strict geometric conditions hold eventually in the
representation coupling.  A sufficiently slowly diverging
deterministic \(j_0(m)\) would then leave
\(\mathfrak j_m>j_0(m)\) with positive probability on that
part, contradicting the bilateral bound just proved.

Proposition~\ref{prop:readout-extension} applies to this
intrinsic exact-record output.  Proposition~\ref{prop:sphere-support}
then transfers, on activation, the nullity of the event that the
extracted threshold equals \(+\infty\) to strict ordinary sphere
interiors, using the common conditional kernel and the actual
field and traversal order.  Equivalently this is the
support-transfer use of \cite[Theorem~7.13 and Lemma~18.1]{P}.
No microscopic circuit event is transferred at this stage.
In a representation coupling, convergence of the compactified
threshold to a finite value makes \(\mathfrak j_m\) eventually
bounded.  The definition of \(\mathfrak j_m\) gives
\eqref{eq:block-approximation}, even with \(C=1\), for that
entry.  Countable intersection gives the simultaneous atlas
statement.

Finally take the two-rooting extraction in
Proposition~\ref{prop:geometry}.  Its ordinary off-root
neighborhoods, with the separately chosen auxiliary mark
pairs, cover the sphere.  Choose finitely many compact sets
\(K_s\) whose interiors cover the sphere and whose closures
have room inside those neighborhoods.  For the root assigned
to \(K_s\), its full time preimage is a compact subset of the
interiors of finitely many used blocks.  Cover that preimage
by interiors of finitely many relative rational intervals whose
closures stay in individual used interiors.  Compactness gives
a positive time margin for this cover.  These positive-length
intervals are entries of the atlas just treated.

There is a uniform positive spatial margin such that every
sphere point lies that far inside at least one \(K_s\).  For a
discrete corner occurrence choose such an \(s\) at its physical
position, and express the occurrence in that set's rooted order.
Uniform tour tracking puts its limiting tour position in \(K_s\)
for all large \(n\), simultaneously for all occurrences.
The time-cover margin and integer rounding then put its cut in
one of the chosen intervals.  Thus every occurrence is covered,
including repeated occurrences of a vertex.  The finite
maximum of their eventually bounded thresholds is finite.
This proves the final assertion of
Theorem~\ref{thm:block-approximation}.
